\section{Interférométrie}

\subsection{Description du dispositif expérimental}

Notre montage utilise un interféromètre de Michelson pour mesurer des déplacements nanométriques. Un laser He-Ne (5 mW, λ = 633 nm) émet un faisceau divisé en deux bras perpendiculaires par une lame séparatrice (en bleu sur la Fig. \ref{fig:schema_inter}). Un miroir fixe et un miroir mobile monté sur l'actuateur piézoélectrique réfléchissent la lumière vers le détecteur. La séparation et "recombinaison" du faisceai crée ainsi des interférences (entre faisceaux vert et rouge sur Fig. \ref{fig:schema_inter}). 
\\\\
En effet, lorsque le piézoélectrique se déplace sous l’effet d’une tension électrique le miroir mobile se translate, car il est solidaire avec l'actuateur. Cela modifie la différence de marche entre les deux bras. Les faisceaux réfléchis interfèrent au niveau de la lame séparatrice. Cela permet de produire sur le détecteur CCD un motif de franges d’interférence (zones claires et sombres alternées) qui sont agrandies par un objectif de microscope (en jaune sur schéma).
\\\\
!!! Décrire plus en détail comment
\begin{figure}[H]
    \centering
    \includegraphics[scale=0.4]{schema_inter.png}
    \label{fig:schema_inter}
    \caption{Schéma du dispositif de l'interféromètre de Michelson.}
\end{figure}

\subsection{Théorie}

L’interféromètre de Michelson permet de relier le déplacement $d$ du miroir mobile à la différence de marche $\delta$ entre les deux bras. Dans notre montage, où l’angle d’incidence est quasi nul, cette relation se simplifie en :
\begin{equation}
    \delta \approx 2d
    \label{eq:delta}
\end{equation}

Un décalage d’une frange correspond à une variation de $\delta$ de $\lambda$ (633 nm). Ainsi, le déplacement du piézoélectrique pour un décalage d’une frange est :
\begin{equation}
    d = \frac{\lambda}{2} = \frac{633 \text{ nm}}{2} = 316,5 \text{ nm}
    \label{eq:deplacement_frange}
\end{equation}

Le coefficient piézoélectrique $d_{33}$ relie la tension appliquée $V$ au déplacement $d$, est défini par :
\begin{equation}
    d = d_{33} \cdot V
    \label{eq:d33_def}
\end{equation}

En mesurant le nombre de franges $n$ décalées pour une tension $\Delta V$, on peut exprimer $d_{33}$ comme :
\begin{equation}
    d_{33} = \frac{n \cdot \frac{\lambda}{2}}{\Delta V}
    \label{eq:d33_calc}
\end{equation}



\subsection{Résultats et Analyse}

\subsubsection{Franges et nano-déplacement}

Nous observons sur le détecteur CCD un motif d'interférence caractéristique composé de franges claires et sombres alternées sur Fig. \ref{fig:franges}.

\begin{figure}[H]
    \centering
    \includegraphics[scale=0.5]{1.jpg}
    \label{fig:franges}
    \caption{Figure d'interférence observée sur le détecteur CCD.}
\end{figure}

Ces franges d'interférence résultent de l'interférence entre les deux faisceaux lumineux réfléchis par les miroirs fixe et mobile. Lorsque ces faisceaux se rejoignent au niveau de la lame séparatrice, leurs ondes lumineuses se superposent. Selon leur différence de phase, la superposition peut être :
Soit constructive quand les crêtes des ondes coïncident (produisant une frange claire), soit destructive quand les crêtes d'une onde coïncident avec les creux de l'autre (produit une frange sombre. Ce phénomène ondulatoire est visible sur Fig. \ref{fig:yay}.

\begin{figure}[H]
    \centering
    \includegraphics[scale=0.5]{yay.png}
    \label{fig:yay}
    \caption{Figure d'interférence observée sur le détecteur CCD.}
\end{figure}

L'espacement entre les franges dépend principalement de la longueur d'onde $\lambda$ du laser et de l'angle des miroirs. En soi, plus $\lambda$ est grande, plus les franges sont espacées. Dans notre montage, les miroirs sont alignés de manière à ce que les faisceaux soient quasi perpendiculaires à la lame séparatrice afin de simplifier la relation entre le déplacement du miroir et le décalage des franges.
\\\\
Ainsi, quand l'angle d'incidence $\theta$ est quasi nul, la différence de marche $\delta$ entre les deux bras de l'interféromètre est donnée par :
\[
\delta \approx 2d
\]
où $d$ est le déplacement du miroir mobile. Ainsi, un décalage d'une frange correspond à une variation de $\delta$ de $\lambda$, soit :
\[
2d = \lambda \implies d = \frac{\lambda}{2} = \frac{633 \text{ nm}}{2} = 316,5 \text{ nm}
\]


%%% Todo
%   inclure voltage puis ca puis coeff
%%%

On peut se demander dès lors pourquoi le déplacement des franges est latéral (perpendiculaire aux franges).
Bien que le piézoélectrique se déplace \textbf{parallèlement à l'axe du CCD}, les franges se déplacent \textbf{perpendiculairement} sur le détecteur. Ce comportement contre-intuitif s'explique par la géométrie de l'interféromètre de Michelson :
\begin{itemize}
    \item Le déplacement du miroir mobile modifie la \textbf{différence de marche $\delta$} entre les deux bras.
    \item Cette modification de $\delta$ entraîne un \textbf{décalage du motif d'interférence} dans une direction orthogonale à l'axe optique principal.
    \item Le détecteur CCD capte ce décalage sous forme d'un mouvement latéral des franges.
\end{itemize}
C'est une propriété fondamentale des interféromètres de Michelson, où le déplacement du miroir se traduit par un décalage latéral des franges sur le détecteur.

\textbf{Calcul du déplacement pour une frange :}
D'après les équations précédentes, un décalage d'\textbf{une frange} sur le détecteur CCD correspond à un déplacement du piézoélectrique de :
\[
d = \frac{\lambda}{2} = \frac{633 \text{ nm}}{2} = 316,5 \text{ nm}
\]

\textbf{Déplacement minimum observable :}
La résolution de notre montage dépend de notre capacité à distinguer un décalage partiel des franges. Si nous pouvons détecter un décalage de \textbf{1/10ème de frange}, alors le déplacement minimum observable du piézoélectrique est :
\[
d_{\text{min}} = \frac{316,5 \text{ nm}}{10} = 31,65 \text{ nm}
\]
Cela démontre la \textbf{précision nanométrique} de la méthode interférométrique, qui permet de mesurer des déplacements extrêmement faibles.

\subsubsection{Coefficient piézoélectrique}

Pour déterminer le coefficient piézoélectrique $d_{33}$, nous avons appliqué des tensions croissantes $V_1, V_2, V_3, \dots$ au piézoélectrique et mesuré le nombre de franges décalées pour chaque tension. Les résultats sont présentés dans le Tableau \ref{tab:deplacement_franges}.

\begin{table}[h]
\centering
\caption{Déplacement du piézoélectrique en fonction de la tension appliquée.}
\begin{tabular}{|c|c|c|}
\hline
\textbf{Tension (V)} & \textbf{Nombre de franges décalées ($n$)} & \textbf{Déplacement $d = n \cdot 316,5$ nm} \\
\hline
$V_1$ & 1 & 316,5 nm \\
\hline
$V_2$ & 2 & 633,0 nm \\
\hline
$V_3$ & 3 & 949,5 nm \\
\hline
$V_4$ & 4 & 1266,0 nm \\
\hline
\end{tabular}
\label{tab:deplacement_franges}
\end{table}

À partir de ces mesures, nous avons tracé la courbe $d(V)$ (voir Figure \ref{fig:superposition}), où $d$ est le déplacement calculé à partir du nombre de franges décalées. La pente de cette courbe donne directement le coefficient piézoélectrique $d_{33}$, défini par :
\[
d_{33} = \frac{\Delta d}{\Delta V} = \frac{n \cdot \frac{\lambda}{2}}{\Delta V}
\]

\begin{figure}[H]
    \centering
    \includegraphics[scale=0.5]{superposition_5_to_20_20.png}
    \label{fig:superposition}
    \caption{Superposition des figures d'interférence pour différentes tensions appliquées (5 V, 10 V, 15 V, 20 V). Chaque courbe montre un décalage progressif des franges avec l'augmentation de la tension.}
\end{figure}

La Figure \ref{fig:coeff} présente la courbe $d(V)$ obtenue expérimentalement. En ajustant une droite à ces données, nous avons déterminé la pente, qui correspond au coefficient piézoélectrique $d_{33}$ :
\[
d_{33} = \text{[valeur mesurée]} \pm \text{[incertitude]} \text{ nm/V}
\]

\begin{figure}[H]
    \centering
    \includegraphics[scale=0.5]{coeff.png}
    \label{fig:coeff}
    \caption{Courbe $d(V)$ permettant de déterminer le coefficient piézoélectrique $d_{33}$. La pente de la droite correspond à la valeur expérimentale de $d_{33}$.}
\end{figure}

\textbf{Comparaison avec la valeur théorique :}
Le fabricant du piézoélectrique (modèle P-117-00) indique un déplacement maximal de 82 µm pour une tension de 150 V, ce qui correspond à un coefficient théorique de :
\[
d_{33,\text{théorique}} = \frac{82 \text{ µm}}{150 \text{ V}} = 546,67 \text{ nm/V}
\]
Les écarts éventuels entre la valeur mesurée et la valeur théorique peuvent être attribués à :
\begin{itemize}
    \item Les incertitudes expérimentales sur la mesure du nombre de franges.
    \item Les imperfections dans l'alignement de l'interféromètre.
    \item Les non-linéarités du matériau piézoélectrique, en particulier à haute tension.
\end{itemize}

En conclusion, cette expérience a permis de valider la précision de la méthode interférométrique pour mesurer des déplacements nanométriques, tout en déterminant expérimentalement le coefficient piézoélectrique du matériau étudié.
